\documentclass[12pt,a4paper]{article}

% Paquetes
\usepackage[utf8]{inputenc}
\usepackage[spanish]{babel}
\usepackage{geometry}
\geometry{margin=2.5cm}
\usepackage{hyperref}
\usepackage{xcolor}
\usepackage{graphicx}
\usepackage{fancyhdr}
\usepackage{titlesec}
\usepackage{enumitem}
\usepackage{booktabs}
\usepackage{array}

% Colores personalizados
\definecolor{azulgit}{RGB}{24,80,50}
\definecolor{grisclaro}{RGB}{245,245,245}

% Configuración de hipervínculos
\hypersetup{
    colorlinks=true,
    linkcolor=azulgit,
    urlcolor=azulgit,
    citecolor=azulgit
}

% Encabezado y pie de página
\pagestyle{fancy}
\fancyhf{}
\fancyhead[L]{\textcolor{azulgit}{\small Programadores Por La Paz}}
\fancyhead[R]{\textcolor{azulgit}{\small Semana 2}}
\fancyfoot[C]{\thepage}
\renewcommand{\headrulewidth}{0.4pt}

% Títulos
\titleformat{\section}
  {\normalfont\Large\bfseries\color{azulgit}}
  {\thesection}{1em}{}
\titleformat{\subsection}
  {\normalfont\large\bfseries\color{azulgit!80!black}}
  {\thesubsection}{1em}{}

\begin{document}

% Título
\begin{titlepage}
    \centering
    \vspace*{2cm}
    
    {\Huge\bfseries\textcolor{azulgit}{DOCUMENTO DE PRESENTACIÓN}\\[0.5cm]}
    {\Large\textcolor{azulgit!70!black}{Tarea Semana 2}}\\[1cm]
    
    \vspace{1cm}
    
    {\Large Terminal, Carpetas, Archivos y Git}\\[2cm]
    
    \vspace{2cm}
    
    {\large\textbf{Programa:}}\\
    {\large Programadores Por La Paz}\\[1.5cm]
    
    {\large\textbf{Universidad de Cartagena}}\\[2cm]
    
    \vfill
    
    {\large\textbf{Estudiante:}}\\
    {\Large Devin Polanco Romero}\\[1cm]
    
    {\large Junio 2026}
    
\end{titlepage}

% Índice
\tableofcontents
\newpage

% Contenido principal
\section{Información General}

\begin{table}[h]
\centering
\begin{tabular}{|>{\bfseries}l|l|}
\hline
\textbf{Campo} & \textbf{Detalle} \\
\hline
Estudiante & Devin Polanco Romero \\
\hline
Programa & Programadores Por La Paz \\
\hline
Institución & Universidad de Cartagena \\
\hline
Actividad & Tarea Semana 2 \\
\hline
Tema & Terminal, Carpetas, Archivos y Git \\
\hline
Fecha & Junio 2026 \\
\hline
\end{tabular}
\caption{Información de la actividad}
\end{table}

\section{Enlace del Repositorio}

El desarrollo de la actividad fue publicado en el siguiente repositorio de GitHub:

\vspace{0.5cm}

\begin{center}
\large
\url{https://github.com/DevCorpSoftwareFactory/Programadores_Para_La_Paz_Semana_2}
\end{center}

\vspace{0.5cm}

\noindent La carpeta \texttt{semana2} contiene todos los archivos solicitados para la entrega de la actividad.

\section{Objetivo de la Semana}

Aprender los fundamentos de la \textbf{terminal}, \textbf{carpetas}, \textbf{archivos}, \textbf{Git} y \textbf{GitHub} para organizar y publicar el trabajo técnico de manera profesional.

\section{Desarrollo de la Actividad}

\subsection{Preguntas de Selección Múltiple}

Se respondieron cuatro preguntas sobre el uso de la terminal:

\begin{enumerate}[label=\textbf{\arabic*.}]
    \item \textbf{¿Qué es la terminal?} \\
    Respuesta: Ejecutar comandos para interactuar con el sistema.
    
    \item \textbf{¿Qué hace el comando \texttt{pwd}?} \\
    Respuesta: Muestra la ubicación actual dentro del sistema de carpetas.
    
    \item \textbf{¿Qué comando lista archivos y carpetas?} \\
    Respuesta: \texttt{ls}
    
    \item \textbf{¿Qué comando crea una nueva carpeta?} \\
    Respuesta: \texttt{mkdir}
\end{enumerate}

\subsection{Comandos de Terminal Utilizados}

Se practicaron los siguientes comandos fundamentales:

\begin{description}
    \item[pwd] Muestra la ubicación actual (directorio de trabajo).
    \item[ls] Lista archivos y carpetas del directorio actual.
    \item[mkdir] Crea una nueva carpeta.
    \item[cd] Navega entre directorios.
    \item[touch] Crea un archivo vacío.
\end{description}

\subsection{Flujo de Trabajo Practicado}

Se siguió el flujo completo de trabajo con Git:

\begin{enumerate}
    \item Verificar ubicación con \texttt{pwd}
    \item Listar contenido con \texttt{ls}
    \item Crear carpeta con \texttt{mkdir}
    \item Navegar con \texttt{cd}
    \item Crear archivos con \texttt{touch}
    \item Preparar cambios con \texttt{git add}
    \item Guardar cambios con \texttt{git commit}
    \item Subir a GitHub con \texttt{git push}
\end{enumerate}

\subsection{Reflexión}

Se desarrolló una reflexión sobre la importancia de la terminal para los desarrolladores, destacando:

\begin{itemize}
    \item Mayor control y precisión sobre el sistema de archivos
    \item Funciona en cualquier entorno (incluyendo servidores)
    \item Facilita la automatización de tareas con scripts
    \item Es habilidad fundamental para ser productivo y colaborativo
\end{itemize}

\section{Archivos Entregados}

La carpeta \texttt{semana2} contiene los siguientes archivos:

\begin{table}[h]
\centering
\begin{tabular}{|l|l|}
\hline
\textbf{Archivo} & \textbf{Contenido} \\
\hline
preguntas-semana2.txt & Respuestas de selección múltiple \\
\hline
comandos-terminal.txt & Comandos utilizados en la práctica \\
\hline
reflexion-semana2.txt & Reflexión sobre la terminal \\
\hline
\end{tabular}
\caption{Archivos de la actividad}
\end{table}

\section{Conclusiones}

\begin{itemize}
    \item Se comprendió el funcionamiento básico de la terminal y su importancia para los desarrolladores.
    \item Se practicaron comandos fundamentales de navegación y gestión de archivos.
    \item Se entendió el flujo de trabajo completo con Git (add, commit, push).
    \item Se reconoció la terminal como herramienta esencial para la productividad y colaboración en desarrollo de software.
\end{itemize}

\section{Referencias}

\begin{itemize}
    \item Documentación oficial de Git: \url{https://git-scm.com/doc}
    \item Documentación de GitHub: \url{https://docs.github.com}
    \item Libro Pro Git: \url{https://git-scm.com/book/es/v2}
\end{itemize}

\end{document}
